\documentclass[twoside,12pt]{book}

\usepackage[utf8]{inputenc}
\usepackage[T1]{fontenc}
\usepackage{mathpazo}
\usepackage{amsmath,amssymb,amsthm}
\usepackage{graphicx}
\usepackage{listings}
\usepackage{xcolor}
\usepackage{enumitem}
\usepackage{tikz}
\usepackage{algorithm}
\usepackage{algpseudocode}
\usepackage{booktabs}
\usepackage{geometry}
\usepackage{fancyhdr}
\usepackage{titlesec}
\usepackage[colorlinks=true, linkcolor=blue, citecolor=blue, urlcolor=blue, plainpages=false, pdfpagelabels]{hyperref}

\geometry{
  papersize={6in,9in},
  margin=0.75in,
  top=0.75in,bottom=0.75in,
  inner=0.85in,outer=0.65in,
  headheight=14pt
}

\pagestyle{fancy}
\fancyhf{}
\fancyhead[LE]{\small\textit{5. Tail Inequalities}}
\fancyhead[RO]{\small\textit{Randomized Algorithms}}
\fancyfoot[C]{\thepage}
\renewcommand{\headrulewidth}{0.4pt}

\definecolor{codegray}{rgb}{0.45,0.45,0.45}
\lstdefinestyle{cppstyle}{
  basicstyle=\ttfamily\scriptsize,
  keywordstyle=\bfseries,
  commentstyle=\itshape\color{codegray},
  numbers=left,
  numberstyle=\tiny\color{codegray},
  numbersep=8pt,
  frame=single,
  framerule=0.4pt,
  rulecolor=\color{codegray},
  backgroundcolor=\color{gray!5},
  breaklines=true,
  tabsize=4,
  showstringspaces=false,
  language=C++,
  morekeywords={nullptr,size_t,auto,const,constexpr,
    unordered_map,vector,pair,mt19937,
    uniform_int_distribution,INT_MAX}
}
\lstset{style=cppstyle}

\theoremstyle{plain}
\newtheorem{theorem}{Theorem}[chapter]
\newtheorem{lemma}[theorem]{Lemma}
\newtheorem{corollary}[theorem]{Corollary}
\newtheorem{proposition}[theorem]{Proposition}
\theoremstyle{definition}
\newtheorem{definition}[theorem]{Definition}
\newtheorem{example}[theorem]{Example}
\newtheorem{remark}[theorem]{Remark}
\theoremstyle{remark}
\newtheorem{exercise}[theorem]{Exercise}
\newtheorem{problem}[theorem]{Problem}
\newenvironment{cpplisting}{\begin{center}\small}{\end{center}}

\newcommand{\Exp}{\operatorname{\mathbb{E}}}
\newcommand{\Var}{\operatorname{Var}}
\newcommand{\Cov}{\operatorname{Cov}}

\begin{document}

\chapter*{5 Tail Inequalities}
\addcontentsline{toc}{chapter}{5 Tail Inequalities}
\markboth{5 Tail Inequalities}{Randomized Algorithms}
\label{ch:tail-inequalities}

In the analysis of randomized algorithms, we frequently need to bound the
probability that a random variable deviates significantly from its expected
value. The Markov inequality provides a generic tool for this purpose, but
its bound is often too loose for practical use.

\begin{definition}[With High Probability]
\label{def:whp}
An event $E_n$ (parametrized by the input size~$n$) occurs \emph{with high
probability} (abbreviated w.h.p.) if $\Pr[E_n] \ge 1 - O(n^{-c})$ for some
constant $c > 0$.  Equivalently, the failure probability is polynomially
small in $n$.  In some contexts, ``w.h.p.'' may require
$\Pr[E_n] \ge 1 - 1/n^c$ for any desired constant~$c$, where the hidden
constant in the algorithm's running time may depend on~$c$.
\end{definition}

Note that the threshold ``w.h.p.'' is stronger than ``with constant
probability'' ($\Pr \ge 2/3$).  By repeating an algorithm $O(\log n)$
times and taking a median (or a union bound), constant-probability
guarantees can typically be amplified to w.h.p.\ guarantees.

In this chapter, we develop
sharper \emph{tail inequalities}---Chernoff bounds, Hoeffding bounds, and
their martingale extensions---that exploit independence and structural
properties to yield exponentially small failure probabilities. These bounds
are among the most frequently used tools in the analysis of randomized
algorithms and have far-reaching applications in computer science,
statistics, and combinatorics.

%======================================================================
\section{The Chernoff Bound}
\label{sec:chernoff}
%======================================================================

We begin with the Chernoff bound, perhaps the most widely used tail
inequality in the analysis of randomized algorithms. The bound applies to
the sum of independent Bernoulli random variables and provides
exponentially decaying bounds on the probability that the sum deviates
from its expectation.

Let $X_1, X_2, \ldots, X_n$ be independent Bernoulli random variables,
where $\Pr[X_i = 1] = p_i$ and $\Pr[X_i = 0] = 1 - p_i$. Define
$X = \sum_{i=1}^{n} X_i$ with expectation $\mu = \mathbb{E}[X] =
\sum_{i=1}^{n} p_i$. Our goal is to bound $\Pr[X \geq (1+\delta)\mu]$
and $\Pr[X \leq (1-\delta)\mu]$ for $\delta > 0$.

The key idea behind the Chernoff bound is the use of the \emph{moment
generating function} (MGF). For any $\lambda > 0$,
\[
\Pr[X \geq a] = \Pr[e^{\lambda X} \geq e^{\lambda a}]
\leq \frac{\mathbb{E}[e^{\lambda X}]}{e^{\lambda a}}
\]
by Markov's inequality. The challenge is to bound
$\mathbb{E}[e^{\lambda X}]$.

\begin{lemma}[MGF Bound for Bernoulli Sums]
\label{lem:mgf}
Let $X_1, X_2, \ldots, X_n$ be independent Bernoulli random variables with
$\mathbb{E}[X_i] = p_i$ and $\mu = \sum_{i=1}^{n} p_i$. Then for all
$\lambda > 0$,
\[
\mathbb{E}[e^{\lambda X}] \leq e^{\mu(e^{\lambda} - 1)}.
\]
\end{lemma}

\begin{proof}
Since the $X_i$ are independent,
\[
\mathbb{E}[e^{\lambda X}] = \prod_{i=1}^{n} \mathbb{E}[e^{\lambda X_i}]
= \prod_{i=1}^{n} \bigl(p_i e^{\lambda} + (1-p_i)\bigr)
= \prod_{i=1}^{n} \bigl(1 + p_i(e^{\lambda} - 1)\bigr).
\]
Using the inequality $1 + x \leq e^x$ (valid for all $x$), each factor
satisfies
\[
1 + p_i(e^{\lambda} - 1) \leq \exp\bigl(p_i(e^{\lambda} - 1)\bigr).
\]
Taking the product,
\[
\mathbb{E}[e^{\lambda X}] \leq \exp\!\Bigl(\sum_{i=1}^{n}
p_i(e^{\lambda} - 1)\Bigr) = e^{\mu(e^{\lambda} - 1)}. \qedhere
\]
\end{proof}

Combining Lemma~\ref{lem:mgf} with Markov's inequality yields the
Chernoff bound in its standard multiplicative form.

\begin{theorem}[Chernoff Bound---Multiplicative Form]
\label{thm:chernoff}
Let $X_1, X_2, \ldots, X_n$ be independent Bernoulli random variables
with $\mathbb{E}[X] = \mu = \sum_{i=1}^{n} p_i$ where $X =
\sum_{i=1}^{n} X_i$. Then for any $\delta > 0$:
\begin{enumerate}[label=(\roman*)]
\item \textbf{Upper tail:}
$\Pr[X \geq (1+\delta)\mu] \leq
   \left(\frac{e^{\delta}}{(1+\delta)^{(1+\delta)}}\right)^{\!\mu}.$
In particular, for $0 < \delta \leq 1$:
\[
\Pr[X \geq (1+\delta)\mu] \leq e^{-\mu\delta^2/3}.
\]
\item \textbf{Lower tail:}
$\Pr[X \leq (1-\delta)\mu] \leq
   \left(\frac{e^{-\delta}}{(1-\delta)^{(1-\delta)}}\right)^{\!\mu}.$
In particular, for $0 < \delta \leq 1$:
\[
\Pr[X \leq (1-\delta)\mu] \leq e^{-\mu\delta^2/2}.
\]
\end{enumerate}
\end{theorem}

\begin{proof}
We prove the upper tail; the lower tail is analogous.

By Markov's inequality and Lemma~\ref{lem:mgf}, for any $\lambda > 0$:
\[
\Pr[X \geq (1+\delta)\mu] \leq
e^{-\lambda(1+\delta)\mu} \cdot e^{\mu(e^{\lambda}-1)}.
\]
Setting $\lambda = \ln(1+\delta) > 0$:
\[
\Pr[X \geq (1+\delta)\mu] \leq
(1+\delta)^{-(1+\delta)\mu} \cdot e^{\mu\delta} =
\left(\frac{e^{\delta}}{(1+\delta)^{(1+\delta)}}\right)^{\!\mu}.
\]
For $0 < \delta \leq 1$, the concavity of $\ln$ and a Taylor expansion
show that $\ln(1+\delta) \geq \delta - \delta^2/2$, which yields
$\delta - (1+\delta)\ln(1+\delta) \leq -\delta^2/3$, giving the simplified
bound $e^{-\mu\delta^2/3}$.

For the lower tail, we choose $\lambda = -\ln(1-\delta) > 0$ and obtain
the analogous bound $e^{-\mu\delta^2/2}$.
\end{proof}

\begin{remark}
The asymmetry in the exponents ($-1/3$ vs.\ $-1/2$) reflects the fact
that the lower tail decays faster than the upper tail. This is a
consequence of the convexity of $e^x$ and the different optimal choices
of $\lambda$.
\end{remark}

\begin{example}[Application to Lower Bounds]
\label{ex:chernoff-lower}
Chernoff bounds provide a simple route to proving lower bounds on the
running time of randomized algorithms. Consider the problem of estimating
the fraction of satisfying assignments of a DNF formula. Each clause in
the DNF is satisfied by a fraction $p_j$ of random assignments. Sampling
$k$ random assignments and checking each against the DNF, the number of
satisfied samples $S$ satisfies $\mathbb{E}[S] = kp$ where $p$ is the
true fraction. By the Chernoff bound, to estimate $p$ within relative
error $\epsilon$, we need $k = O(\ln(1/\delta)/(p\epsilon^2))$ samples.
This immediately implies a lower bound: no algorithm using fewer samples
can achieve relative error $\epsilon$ with probability exceeding $1 -
\delta$.
\end{example}

\subsection{Derandomization via Conditional Expectations}

Chernoff bounds also serve as the foundation for \emph{derandomization}
of randomized algorithms. The method of conditional expectations allows
us to systematically replace random choices with deterministic ones while
preserving the expectation bound. The key observation is that
$\mathbb{E}[X \mid X_1 = x_1, \ldots, X_k = x_k]$ is itself a sum of
independent Bernoulli random variables, to which the Chernoff bound
applies. By evaluating conditional expectations, one can deterministically
construct a solution that is at least as good as the expectation.

%======================================================================
\section{Routing in a Parallel Computer}
\label{sec:routing}
%======================================================================

One of the classic applications of Chernoff bounds in computer science is
to the analysis of randomized routing algorithms in parallel architectures.

Consider a parallel computer with $n$ processors connected to $n$ memory
modules. At each time step, each processor independently selects a memory
module uniformly at random and sends a packet to that module. If multiple
packets are destined for the same module, they queue up and are served
one at a time. The routing is completed when all packets have been
delivered, and the \emph{congestion} is the maximum queue length at any
module.

The expected number of packets arriving at any given memory module is
$\mu = n \cdot (1/n) = 1$. Thus, we expect each module to receive about
one packet. However, the interesting question is whether the maximum
congestion is close to this expectation.

\begin{theorem}[Routing Congestion]
\label{thm:routing}
In the random routing model described above, the maximum queue length is
at most $\frac{\ln n}{\ln \ln n}$ with probability at least $1 - n^{-c}$
for any fixed constant $c > 0$.
\end{theorem}

\begin{proof}
Let $q_j$ denote the queue length at module $j$. Since each of the $n$
processors independently chooses module $j$ with probability $1/n$, we
have $q_j = \sum_{i=1}^{n} X_{ij}$ where $X_{ij}$ is the indicator
that processor $i$ sends to module $j$. Thus $q_j$ is a sum of $n$
independent Bernoulli random variables with $p_i = 1/n$ for each, giving
$\mu = 1$.

By the Chernoff bound (Theorem~\ref{thm:chernoff}(i)), for any $t > 0$:
\[
\Pr[q_j \geq t] \leq \left(\frac{e^{t-1}}{t^t}\right).
\]
Setting $t = \frac{\ln n}{\ln \ln n}$ and applying a union bound over
all $n$ modules:
\[
\Pr\!\Bigl[\max_j q_j \geq t\Bigr]
\leq n \cdot \Pr[q_1 \geq t]
\leq n \cdot \left(\frac{e}{t}\right)^t
= n \cdot \exp\bigl(t(1 - \ln t)\bigr).
\]
Substituting $t = \frac{\ln n}{\ln \ln n}$:
\[
t(1 - \ln t) = \frac{\ln n}{\ln\ln n}
\left(1 - \ln\!\left(\frac{\ln n}{\ln\ln n}\right)\right)
= \frac{\ln n}{\ln\ln n}\left(1 - \ln\ln n + \ln\ln\ln n\right),
\]
which is $O(-\ln n)$ for large $n$. Hence the failure probability is
$n \cdot e^{-\Omega(\ln n)} = n^{-c}$ for suitable $c > 1$.
\end{proof}

\begin{remark}
The bound of $\frac{\ln n}{\ln \ln n}$ is asymptotically tight for this
model: it can be shown that with high probability, the maximum congestion
is at least $\Omega\!\left(\frac{\ln n}{\ln \ln n}\right)$. The key
insight is that the union bound, combined with the exponential decay of
the Chernoff bound, converts the per-module exponential tail into a
high-probability global guarantee.
\end{remark}

The following C++ code illustrates this phenomenon experimentally.
We simulate random routing for varying numbers of processors and compare
the observed maximum congestion against the theoretical bound.

\begin{cpplisting}
\begin{lstlisting}[caption={Routing simulation (see \texttt{routing.h})}]
int simulate_routing(int n, mt19937& rng) {
    std::vector<int> queues(n, 0);
    uniform_int_distribution<int> dist(0, n - 1);
    for (int i = 0; i < n; i++)
        queues[dist(rng)]++;
    return *max_element(queues.begin(), queues.end());
}
\end{lstlisting}
\end{cpplisting}

%======================================================================
\section{A Wiring Problem}
\label{sec:wiring}
%======================================================================

The Chernoff bound finds another elegant application in the analysis of
wiring problems in VLSI design. Consider the following model: we have $n$
signal sources and $n$ signal sinks placed on a line. Each source $i$ is
connected to a unique sink $\pi(i)$, where $\pi$ is a random permutation.
Each wire is laid out along the line, and the total wire length is
$\sum_{i=1}^{n} |i - \pi(i)|$.

For a random permutation, the displacement $|i - \pi(i)|$ is a random
variable with $\mathbb{E}[|i - \pi(i)|] = \Theta(\sqrt{n})$ for a
typical index $i$ (more precisely, the expected displacement for a
uniformly random position is $\mathbb{E}[|i - \pi(i)|] \leq \sqrt{n}$).
By linearity of expectation, the expected total wire length is
$\Theta(n^{3/2})$.

The Chernoff bound provides a concentration result: with high probability,
the total wire length is close to its expectation. More precisely, let
$W = \sum_{i=1}^{n} |i - \pi(i)|$. Although the terms $|i - \pi(i)|$ are
not independent (they are determined by the single permutation $\pi$),
one can decompose the analysis into independent contributions using the
method of \emph{Hoeffding decomposition} or by bounding the sum by
independent random variables.

\begin{theorem}[Wire Length Concentration]
\label{thm:wiring}
For a random permutation $\pi$ on $\{1, \ldots, n\}$, the total wire
length $W = \sum_{i=1}^{n} |i - \pi(i)|$ satisfies, for any $\delta > 0$:
\[
\Pr\!\bigl[|W - \mathbb{E}[W]| \geq \delta \cdot \mathbb{E}[W]\bigr]
\leq 2\exp\!\bigl(-\Omega(\delta^2 n)\bigr).
\]
\end{theorem}

This result has practical implications: it guarantees that random wiring
layouts in VLSI circuits will, with high probability, have wire lengths
close to the expected value, enabling reliable performance prediction.

%======================================================================
\section{Martingales}
\label{sec:martingales}
%======================================================================

The Chernoff bound requires independence of the underlying random
variables, which limits its applicability. Many natural random variables
exhibit sequential dependence---for example, the number of edges in a
random graph, the height of a random search tree, or the score in a
sequential process. \emph{Martingales} provide the right framework for
analyzing such quantities.

\begin{definition}[Martingale]
\label{def:martingale}
A sequence of random variables $X_0, X_1, X_2, \ldots, X_n$ is called a
\emph{martingale} with respect to a filtration (sequence of
$\sigma$-algebras) $\mathcal{F}_0 \subseteq \mathcal{F}_1 \subseteq
\cdots \subseteq \mathcal{F}_n$ if:
\begin{enumerate}[label=(\roman*)]
\item $X_i$ is $\mathcal{F}_i$-measurable for each $i$;
\item $\mathbb{E}[|X_i|] < \infty$ for each $i$;
\item $\mathbb{E}[X_i \mid \mathcal{F}_{i-1}] = X_{i-1}$ for each
      $i \geq 1$.
\end{enumerate}
Intuitively, a martingale models a ``fair game'': the expected future
value, given the present, equals the present value.
\end{definition}

\begin{example}[Random Graph Edge Count]
\label{ex:graph-martingale}
Consider the Erd\H{o}s--R\'{e}nyi random graph $G(n, p)$, and let $X_k$
denote the number of edges in the subgraph induced by the first $k$
vertices. Then $X_0, X_1, \ldots, X_n$ forms a martingale (with respect
to the natural filtration). Indeed, when the $(k+1)$-th vertex is added,
each potential edge to the existing $k$ vertices is included independently
with probability $p$. The conditional expectation of the new edges is
$kp$, so $\mathbb{E}[X_{k+1} \mid X_k] = X_k + kp$ --- wait, this is
not a martingale but a \emph{submartingale}. A proper martingale arises
by considering $Z_k = X_k - \mathbb{E}[X_k] = X_k - \binom{k}{2}p$.
\end{example}

The power of martingale theory lies in concentration inequalities that
generalize the Chernoff bound.

\begin{theorem}[Azuma--Hoeffding Inequality]
\label{thm:azuma}
Let $X_0, X_1, \ldots, X_n$ be a martingale with bounded differences:
\[
|X_i - X_{i-1}| \leq c_i \quad \text{for all } i = 1, \ldots, n,
\]
where $c_1, c_2, \ldots, c_n$ are constants. Then for any $t > 0$:
\[
\Pr\!\bigl[|X_n - X_0| \geq t\bigr]
\leq 2\exp\!\left(-\frac{t^2}{2\sum_{i=1}^{n} c_i^2}\right).
\]
\end{theorem}

\begin{proof}
We prove the upper tail $\Pr[X_n - X_0 \geq t]$; the lower tail follows
by symmetry. Define the \emph{Doob martingale} differences $D_i = X_i -
X_{i-1}$, so $X_n - X_0 = \sum_{i=1}^{n} D_i$.

For any $\lambda > 0$:
\[
\Pr[X_n - X_0 \geq t] = \Pr\!\left[e^{\lambda \sum D_i}
\geq e^{\lambda t}\right]
\leq e^{-\lambda t} \cdot \mathbb{E}\!\left[e^{\lambda \sum D_i}\right].
\]
Since $\mathbb{E}[D_i \mid \mathcal{F}_{i-1}] = 0$ (by the martingale
property) and $|D_i| \leq c_i$, the Hoeffding lemma gives:
\[
\mathbb{E}\!\left[e^{\lambda D_i} \mid \mathcal{F}_{i-1}\right]
\leq e^{\lambda^2 c_i^2 / 2}.
\]
Taking iterated expectations:
\[
\mathbb{E}\!\left[e^{\lambda \sum D_i}\right]
\leq \exp\!\left(\frac{\lambda^2}{2}\sum_{i=1}^{n} c_i^2\right).
\]
Minimizing over $\lambda > 0$ by setting $\lambda = t / \sum c_i^2$
yields the claimed bound.
\end{proof}

\begin{corollary}[Concentration for Bounded-Difference Functions]
\label{cor:bounded-diff}
Let $f(x_1, \ldots, x_n)$ be a function satisfying the
\emph{bounded differences} condition: changing any single input $x_i$
changes $f$ by at most $c_i$, while leaving the other inputs fixed.
If $X_1, \ldots, X_n$ are independent random variables, then
$Z = f(X_1, \ldots, X_n)$ satisfies:
\[
\Pr\!\bigl[|Z - \mathbb{E}[Z]| \geq t\bigr]
\leq 2\exp\!\left(-\frac{t^2}{2\sum_{i=1}^{n} c_i^2}\right).
\]
\end{corollary}

\begin{proof}
Define the Doob martingale $Z_k = \mathbb{E}[Z \mid X_1, \ldots, X_k]$.
Then $Z_0 = \mathbb{E}[Z]$ and $Z_n = Z$. The bounded differences
condition ensures $|Z_k - Z_{k-1}| \leq c_k$, so Theorem~\ref{thm:azuma}
applies.
\end{proof}

\begin{example}[Random Graph Applications]
\label{ex:graph-azuma}
The Azuma--Hoeffding inequality is particularly powerful for random
graphs. For the chromatic number $\chi(G)$ of $G(n, 1/2)$, we have
$\chi(G) = f(X_1, \ldots, X_m)$ where $X_e$ are the edge indicators and
$f$ satisfies bounded differences with $c_e = 1$ for all $e$. Thus:
\[
\Pr\!\bigl[|\chi(G) - \mathbb{E}[\chi(G)]| \geq t\bigr]
\leq 2\exp\!\left(-\frac{t^2}{m}\right)
\]
where $m = \binom{n}{2}$. This shows that the chromatic number of a
random graph is highly concentrated around its expectation.
\end{example}

\begin{remark}
The Azuma--Hoeffding inequality is tight up to constant factors in the
exponent when all $c_i$ are equal. More refined inequalities (such as the
Bennett and Bernstein inequalities) provide better constants when the
$c_i$ are non-uniform or when additional information about the
distribution is available.
\end{remark}

%======================================================================
\addcontentsline{toc}{section}{Notes}
\section*{Notes}
%======================================================================

The Chernoff bound originated in the work of \textbf{Chernoff} (1952) on
the asymptotic efficiency of tests and estimates. The multiplicative form
presented here was developed and popularized by
\textbf{Hoeffding} (1963), who also established the more general bounds
for sums of independent bounded random variables. An earlier, independent
result by \textbf{Hoeffding} (1956) on the Hoeffding inequality for
bounded random variables predates the Chernoff-style bound in the
literature.

The routing application (Section~\ref{sec:routing}) is based on the
analysis of packet routing in parallel architectures, which has been a
central problem in parallel computing since the 1980s. The bound of
$O(\ln n / \ln \ln n)$ for random routing is a classical result; see
Leighton, Maggs, and Rao (1994) for related results on oblivious routing.

The Azuma--Hoeffding inequality (Theorem~\ref{thm:azuma}) is named after
\textbf{Azuma} (1967), who proved a version for martingales with bounded
differences, and \textbf{Hoeffding} (1963), who proved the independent
case. The elegant connection between bounded differences and martingale
concentration was developed by \textbf{McDiarmid} (1989) in his
foundational paper on the ``method of bounded differences.'' The survey by
\textbf{Spencer} (1993) provides an excellent introduction to these
techniques with numerous applications in combinatorics.

For further reading, we recommend the comprehensive treatment in
\emph{Probability and Computing} by Mitzenmacher and Upfal (2005), and
the monograph \emph{Concentration Inequalities} by Boucheron, Lugosi, and
Massart (2013).

%======================================================================
\addcontentsline{toc}{section}{Problems}
\section*{Problems}
%======================================================================

\begin{problem}
Let $X_1, \ldots, X_n$ be independent Bernoulli random variables with
$\Pr[X_i = 1] = 1/n$ for all $i$, and let $X = \sum X_i$. Use the
Chernoff bound to show that $\Pr[X \geq 2] \leq n/e^2$ for $n \geq 4$.
Compare with the exact value obtained from the Poisson approximation.
\end{problem}

\begin{problem}
Prove that the lower tail bound in Theorem~\ref{thm:chernoff}(ii) is
tight up to constant factors in the exponent. That is, exhibit a sequence
of independent Bernoulli random variables for which $\Pr[X \leq
(1-\delta)\mu] \geq e^{-\Theta(\mu\delta^2)}$.
\end{problem}

\begin{problem}
In the routing problem of Section~\ref{sec:routing}, suppose each
processor sends to a \emph{random} destination chosen uniformly from
$\{1, \ldots, n\}$, but now there are $m$ rounds. After each round, any
delivered packets are removed. Show that the expected total time to deliver
all packets is $O(\ln n)$ rounds with high probability.
\end{problem}

\begin{problem}
Let $G = G(n, 1/2)$ be the Erd\H{o}s--R\'{e}nyi random graph. Using the
Azuma--Hoeffding inequality, prove that the number of triangles in $G$ is
concentrated within $O(n^{3/2})$ of its expectation.
\end{problem}

\begin{problem}
Let $f(x_1, \ldots, x_n) = \max_{1 \leq i \leq n} x_i$ where $x_i \in
[0, 1]$. Show that $f$ satisfies the bounded differences condition with
$c_i = 1$ for all $i$. Apply Corollary~\ref{cor:bounded-diff} to the
maximum of $n$ independent uniform $[0,1]$ random variables and compare
the Azuma--Hoeffding bound with the exact distribution.
\end{problem}

\begin{problem}
Prove that if $X_0, X_1, \ldots, X_n$ is a martingale with $|X_i -
X_{i-1}| \leq 1$ for all $i$, then $\Pr[|X_n - X_0| \geq t] \leq
2e^{-t^2/(2n)}$. This is sometimes called the ``simple'' Azuma
inequality.
\end{problem}

\begin{problem}
Consider a random variable $Z = \prod_{i=1}^{n} Y_i$ where $Y_i$ are
independent, each taking value $2$ with probability $1/2$ and value $1/2$
with probability $1/2$. Taking logarithms and applying the Chernoff bound,
show that $\Pr[Z \geq 2^{n/2 + t\sqrt{n}}] \leq e^{-\Omega(t^2)}$.
Interpret this as a ``multiplicative'' Chernoff bound.
\end{problem}

\begin{problem}
Let $A$ be a random $n \times n$ adjacency matrix of $G(n, p)$ with $p =
c/n$ for a constant $c > 0$. The largest eigenvalue $\lambda_{\max}(A)$ of
$A$ satisfies $|\lambda_{\max}(A) - c| \leq O(\sqrt{c})$ with high
probability. Use the method of bounded differences to prove a concentration
inequality for $\lambda_{\max}(A)$.
\end{problem}

\begin{problem}
Show that the moment generating function bound in Lemma~\ref{lem:mgf}
gives a tighter result than directly applying Markov's inequality to $X$
without the exponential transformation. Specifically, compare the bounds
obtained for $\Pr[X \geq 2\mu]$ using (a) Markov's inequality on $X$
directly and (b) the Chernoff bound.
\end{problem}

\begin{problem}
\textbf{(Chernoff for Poisson.)} Let $X \sim \text{Poisson}(\lambda)$.
Prove that $\Pr[X \geq (1+\delta)\lambda] \leq
\left(\frac{e^{\delta}}{(1+\delta)^{(1+\delta)}}\right)^{\lambda}$ for
$\delta > 0$, using the same MGF technique as in the proof of
Theorem~\ref{thm:chernoff}.
\end{problem}

\end{document}
